% Abhakseite „Das kann ich“ – Zone nur im Gesamt (\mitzone)
%% TODO Ich-kann-Sätze: Platzhalter wie in den Aufgabendateien
\begin{abhakseite}
\ifdefined\mitzone
\abhakgruppe{Kennst du schon}
\abhak{1}{Dezimalzahlen und ganze Zahlen mal und geteilt durch eine Zahl (auch nicht durch Zehnerpotenzen: dreißig geteilt durch fünfzig), Kommaverschiebung bei zehn, hundert, tausend}
\abhak{2}{Längeneinheiten umrechnen (m in cm und zurück, km in m, mm in cm), Ergebnis in einer sinnvollen Einheit angeben}
\abhak{3}{Vielfache und Teiler, „doppelt“, „halb“, „dreifach“, „ein Drittel“ als Rechnung}
\abhak{4}{Dreisatz und fester Faktor („pro Portion“; k = Preis je Stück)}
\abhak{5}{Strecken auf Millimeter messen und zeichnen, Rechteck mit dem Geodreieck, Kreis mit dem Zirkel (Radius als halber Durchmesser), Parallelen mit dem Geodreieck zeichnen und erkennen}
\abhak{6}{Draufsicht, Netz und Schrägbild eines Körpers unterscheiden; Draufsicht eines Zylinders ist ein Kreis, eines Quaders ein Rechteck}
\abhak{7}{Verhältnis lesen und als Division schreiben („drei zu fünf“ als Bruch und als Quotient, auch als Dezimalzahl)}
\abhak{8}{Gleichung mit einer Variablen und Verhältnisgleichung nach x umstellen (x geteilt durch die bekannte Strecke gleich dem Verhältnis der zwei anderen; auf beiden Seiten mal nehmen)}
\abhak{9}{Winkel messen und vergleichen, Winkelsumme im Dreieck (zwei Winkel gleich → der dritte auch)}
\abhak{10}{Längeneinheiten umrechnen (m in cm und zurück, km in m, mm in cm), Ergebnis in einer sinnvollen Einheit angeben – Fehler finden}
\abhak{11}{Längeneinheiten umrechnen (m in cm und zurück, km in m, mm in cm), Ergebnis in einer sinnvollen Einheit angeben – selbst rechnen}
\fi
\abhakgruppe{Einheit 1 · Maßstab}
\abhak{12}{Gleiche Einheit?}
\abhak{13}{Maßstab umrechnen}
\abhak{14}{Maßstabsgerecht zeichnen}
\abhak{15}{Maßstab umrechnen – Fehler finden, Begründen, Anwendung, Darstellungswechsel}
\abhakgruppe{Einheit 2 · Zentrische Streckung und Ähnlichkeit}
\abhak{16}{Zentrische Streckung und Ähnlichkeit}
\abhak{17}{Zentrum aus Original und Bild finden}
\abhak{18}{Zentrische Streckung und Ähnlichkeit – Fehler finden, Begründen, Anwendung}
\abhakgruppe{Einheit 3 · Strahlensätze}
\abhak{19}{Strahlensätze}
\abhak{20}{Voraussetzung prüfen}
\abhak{21}{Strecke in n gleiche Teile teilen}
\abhak{22}{Strahlensätze – Fehler finden, Begründen, Anwendung, Darstellungswechsel}
\end{abhakseite}
